\documentclass[11pt, a4paper]{article}

\usepackage[T1]{fontenc}
\usepackage[utf8]{inputenc}
\usepackage{tgpagella}
\usepackage[spanish, es-noshorthands]{babel}
\usepackage{amsmath, amssymb, bm}
\usepackage{tabularx}
\usepackage{booktabs}
\usepackage[most]{tcolorbox}
\usepackage[margin=2.5cm]{geometry}
\usepackage{ragged2e}
\usepackage{fancyhdr}

\pagestyle{fancy}
\fancyhf{}
\setlength{\headheight}{16pt}
\lhead{\textbf{UCB Sede Tarija} | Ing. Mecatrónica}
\chead{}
\rhead{IMT-121 Circuitos Electrónicos I | Remedial Lab A}
\lfoot{Prof: M.Sc. Ing. Bernardo Quiroga Turdera}
\rfoot{Página \thepage}
\renewcommand{\headrulewidth}{0.4pt}

\definecolor{ucbblue}{RGB}{0, 51, 102}

\begin{document}

\begin{tcolorbox}[colback=red!5!white, colframe=red!70!black, title=\textbf{Triage Residual: Cierre Obligatorio de Experiencia A}]
    \centering \textbf{Formulario de Remediación Acotada}
\end{tcolorbox}

\noindent \textit{Aviso: Si su equipo no ha logrado consolidar los parámetros base de su red, NO debe iniciar el montaje de la Experiencia B. Identifique su rezago a continuación y resuélvalo en los próximos 15 minutos.}

\vspace{0.4cm}
\textbf{Equipo:} \hrulefill \quad \textbf{Prioridad Residual Única Asignada:} \hrulefill

\vspace{0.5cm}
\section*{Checklist de Evidencia Mínima}
Marque únicamente lo que \textbf{ya está validado y documentado} con datos en su hoja de registro. Lo que quede desmarcado es su prioridad inmediata.

\begin{itemize}\setlength\itemsep{0.3em}
    \item[$\square$] \textbf{Topología Funcional:} Circuito base montado, tierras unificadas, sin cortocircuitos indeseados.
    \item[$\square$] \textbf{Superposición Experimental Completa:} Mediciones de $V_a^{(V1)}$ y $V_a^{(V2)}$ registradas (desactivando las fuentes con puentes a tierra, no abriendo el circuito). La suma algebraica coincide con la lectura total.
    \item[$\square$] \textbf{Voltaje de Circuito Abierto:} $V_{oc} = V_{th}$ medido correctamente retirando la carga del puerto.
    \item[$\square$] \textbf{Resistencia Equivalente:} $R_{th}$ medida con la red totalmente desenergizada, mirando desde $a-b$.
    \item[$\square$] \textbf{Comparación con SPICE:} Se verificaron las desviaciones relativas entre la simulación y la protoboard.
    \item[$\square$] \textbf{Discrepancia Documentada:} Tienen escrita al menos una hipótesis técnica que explique los errores (ej. tolerancia real de las resistencias $\neq$ nominal).
\end{itemize}

\vspace{1cm}
\begin{center}
    \textbf{Firma del Instructor (Habilitación para iniciar Experiencia B):} \rule{6cm}{0.4pt}
\end{center}

\end{document}
